% npj Digital Medicine: abstract up to 150 words, no subheadings, single
% paragraph. The full-length version is kept in
% 01_abstract_LONG_backup.tex.bak for journals without a hard limit.
\begin{abstract}
\noindent
Medical tasks mapping a low-cost, widely available source modality
onto a costly target are reported as a cohort-average metric that
hides which targets are reachable and why others are not.
We treat them as imperfect translation and report per-target
translatability.
Our three-regime taxonomy---encoder saturation, inter-pattern
co-occurrence, out-of-ontology---is derived as the redundant,
synergistic and unique atoms of a partial information decomposition of
the source ontology against the encoder, paired with a
population-level Conclusive/Inconclusive verdict and the six-level
\lchai{} explanation stack.
On TCGA-LUAD (687 slides, 668 patients) TP53 and EGFR saturate the
encoder, KRAS is synergistic, STK11 and KEAP1 fall outside it, and
RBM10 is only underpowered.
KRAS forces the levels apart: Inconclusive at 0.609 cohort-wide, yet
0.965 on one slide.
Needing only an expert source ontology and a learned encoder, the
decomposition supports two pre-registered public-data tests: glioma
{MRI} to genotype under {VASARI}, and electrocardiogram to laboratory
state under {PTB-XL}'s graded {SCP-ECG} statements.
\end{abstract}

\paragraph{Keywords.}
Cross-domain medical prediction; translatability; partial information
decomposition; lung adenocarcinoma; mutation prediction; radiogenomics;
explainable AI; multiple-instance learning; fuzzy logic; knowledge
graph; TCGA-LUAD; BraTS; PTB-XL.
